\documentclass[nonewpage]{xourse}


% Start \part op een nieuwe pagina (Ximera doet dat niet per default)
\xpretocmd{\part}{\clearpage}{}{}

% \documentclass{xourse}
%in
%\title{\includegraphics[width=0.7\textwidth]{pic5/elektromagnetisme.jpg}\\[2cm]Fysica 5de jaar - Wetenschappen}
\title{Fysica 5de jaar - Niet wetenschappen}
\author{Daan Lipkens}
\begin{document}
\begin{abstract}
    Cursus Fysica 5de jaar niet wetenschappelijke richting.
\end{abstract}
\maketitle
%HOE MAAK IK DIT IN TITLE BLAD 
%\begin{figure}[h]
% \begin{center}
% \includegraphics[width=0.85\textwidth]{pic5/elektromagnetisme.jpg}
% \caption{De Joint European Torus (JET) is het grootste kernfusie-experiment ter wereld. Met behulp van krachtige elektromagnetische velden houdt deze tokamak het extreem hete plasma op zijn plaats, zodat de wandmaterialen niet smelten.}
% \end{center}
% \end{figure}
%\xmtitle{Fysica 5de jaar - Wetenschappen}{}
% \part voor delen
%\chapterstyle voor hoofdstukken >> \activity{hoofdstuk ....}
%\sectionstyle voor delen in het hoofdstuk >> \activity{met onderdelen}
% Inleiding, informatie, ...


% INFO
% toepassingen onder \begin{example}
% Quizjes als \question met \begin{selectAll}
% voorbeeld oefeningen met align x &= berekening && (uitleg)
% komma's met {,} in math mode

%\frontmatter
\part{Inleiding}
\sectionstyle
\activity{praktische_richtlijnen.tex}
\activity{hoe_fysica_te_studeren_5.tex}

%\mainmatter
%--------ELEKTRICITEIT----------------
\part{Elektrische krachtwerking} %parts nog aanpassen
%COULOMBKRACHT of ELEKTRISCHE KRACHTWERKING
%picture path: pic5/elektriciteit/coulomb/

\chapterstyle

\activity{elektrische_krachtwerking_inleiding.tex} %inleiding hoofdstuk coulombkracht

\sectionstyle

\activity{elektrische_krachtwerking_lading.tex} %ladingen en soorten ladingen
%\activity{elektrische_krachtwerking_tribo-elekrisch_effect.tex} %tribo-elektrisch effect --> extra
\activity{elektrische_krachtwerking_coulombkracht_1uur.tex} %coulombkracht zonder cosinus regel
\activity{elektrische_krachtwerking_materialen_en_geleiding_1uur.tex} % isolator en geleider
\activity{elektrische_krachtwerking_aantrekken_ongeladen_deeltjes_1uur.tex} % polarisatie zonder toepassingen
\activity{elektrische_krachtwerking_elektroscoop.tex}
\activity{elektrische_krachtwerking_opladen_inductie.tex}
\activity{elektrische_krachtwerking_samenvatting.tex}
%oefeningen
\chapterstyle
\activity{elektrische_krachtwerkingoef_inleiding.tex}
\sectionstyle
\activity{elektrische_krachtwerkingoef_vragen.tex}
\activity{elektrische_krachtwerkingoef_basis.tex}
\activity{elektrische_krachtwerkingoef_1lijn.tex}
\activity{elektrische_krachtwerkingoef_driehoek_1uurs.tex}



\part{Elektrisch veld}
%ELEKTRISCH VELD
%picture path: pic5/elektriciteit/veld/
 
\chapterstyle
\activity{elektrischveld_inleiding.tex}
 
 
\sectionstyle
\activity{elektrischveld_puntlading.tex} % E-veld van puntlading
\activity{elektrischveld_veldlijnen.tex} % veldlijnen en soorten velden
\activity{elektrischveld_materie_in_veld.tex} % geleider en isolator in veld + toepassingen
% \activity{elektrischveld_samenvatting.tex} %todo
 
% %oefeningen
% \chapterstyle
% \activity{elektrischveldoef_inleiding.tex}
% \sectionstyle
% \activity{elektrischveldoef_vragen.tex}
% \activity{elektrischveldoef_basis.tex}
% \activity{elektrischveldoef_1lijn.tex}
% \activity{elektrischveldoef_driehoek.tex}



\part{Potentiële elektrische energie bij homogeen veld en spanning}
%POTENTIËLE ELEKTRISCHE ENERGIE BIJ HOMOGEEN VELD EN SPANNING
%picture path: /pic5/elektriciteit/energie/
% \chapterstyle

% \sectionstyle

%--------MAGNETISME----------------
\part{Elektromagnetisme}
\part{Lorentzkracht}
\part{Inductiespanning}

\newpage
%--------KERNFYSICA----------------
\part{Deeltjesfysica}\label{h:deeltjesfysica}
\part{De atoomkern} %deeltjesfysica apart deel?
\part{Radioactiviteit}
\part{Radioactiviteit en de mens}


%---------APPENDIX---------------
%\appendix
\part{Appendix} %nummering anders

\chapterstyle
\activity{vectoren_inleiding.tex}

\sectionstyle
\activity{vectoren_scalar.tex}
\activity{vectoren_component.tex}
\activity{vectoren_som.tex}
\activity{vectoren_verschil.tex}
\activity{vectoren_scalaire_vermenigvuldiging.tex}
\activity{vectoren_vectorieel_product.tex}


\chapterstyle
\activity{standaardvorm_inleiding.tex}\label{hoofdstuk:standaardvorm}

\sectionstyle
\activity{standaardvorm_si.tex}
\activity{standaardvorm_bc.tex}
\activity{standaardvorm_oefeningen.tex}

%Vector rekenen
    %vectoren
    %cosinus regel

%nuclide tabel
%standaard vorm
%logaritme en exponentiële functies
%---------Bibliografie
\part{Bibliografie}

\end{document}
